\documentclass{article}
\usepackage{amsmath, amssymb, amsthm}

\setlength{\parindent}{0pt}

\newtheorem*{claim}{Claim}
\newtheorem{lemma}{Lemma}

\begin{document}

\begin{lemma}
    A matrix $A$ is invertible if and only if its echelon form has a pivot in every column and every row.
\end{lemma}

\begin{lemma}
    If a matrix $A$ is invertible, then $A^T$ is also invertible.
\end{lemma}

\begin{claim}
    If the columns of a square $(n \times n)$ matrix $A$ are linearly independent, so are the rows of $A^3$.
\end{claim}

% i think this proof and problem_3-6's proof are way overcomplicated (because i'm stupid), but i'm too lazy to 
% rewrite both of them AND i think they're still correct so...
\begin{proof}
    Suppose $A^3\mathbf{x} = \mathbf{0}$. We first show that $\mathbf{x} = \mathbf{0}$. Let 
    $\mathbf{y} = A\mathbf{x}$ and $\mathbf{z} = A\mathbf{y}$, so 
    \[
    \mathbf{0} = A^3\mathbf{x} = A(A(A\mathbf{x})) = A(A\mathbf{y}) = A\mathbf{z}
    \]
    Since the columns of $A$ are linearly independent, $A\mathbf{z} = \mathbf{0}$ has only the trivial solution 
    $\mathbf{z} = \mathbf{0}$. Applying linear independence of the columns of $A$ to 
    $\mathbf{z} = A\mathbf{y} = \mathbf{0}$, we get $\mathbf{y} = \mathbf{0}$. Applying linear independence of 
    the columns of $A$ once more to $\mathbf{y} = A\mathbf{x} = \mathbf{0}$, we get $\mathbf{x} = \mathbf{0}$.
    Thus $A^3\mathbf{x} = \mathbf{0}$ has only the trivial solution $\mathbf{x} = \mathbf{0}$, so the columns of 
    $A^3$ are linearly independent.

    \medskip

    It remains to show that the rows of $A^3$ are linearly independent. Firstly, because 
    $A^3\mathbf{x} = \mathbf{0}$ has only the trivial solution (or in other words the trivial solution is unique), 
    there are no free variables in the echelon form of $A^3$, so every column has a pivot. Furthermore, since there 
    are $n$ columns, there are $n$ pivots, and by definition of a pivot there can only be at most one pivot in each 
    row. Since there are $n$ rows, every row has a pivot. Hence by Lemma 1, $A^3$ is invertible.

    \medskip

    Furthermore, by Lemma 2, $(A^3)^T$ is also invertible, and by Lemma 1, $(A^3)^T$ has a pivot in every column 
    and every row in its echelon form. Since there is a pivot in every column, there are no free variables, so 
    $(A^3)^T\mathbf{v} = \mathbf{0}$ has only the trivial solution $\mathbf{v} = \mathbf{0}$, so the columns of 
    $(A^3)^T$ are linearly independent. Since the columns of $(A^3)^T$ are linearly independent, the rows of 
    $A^3$ are linearly independent by definition of transpose.

    \medskip

    Hence if the columns of a square matrix $A$ are linearly independent, so are the rows of $A^3$.
\end{proof}

\end{document}
